\chapter{Метод}
\label{ch:method}

\section{Суррогатная функция}
\label{sec:surrogate}

Прямая оптимизация задачи~\eqref{eq:moe_likelihood} требует многократного
обучения нейронных сетей для множества кандидатов-архитектур, что
непомерно дорого в случае MoE, где каждая оценка уже включает $K$ сетей.
Чтобы преодолеть это узкое место, точное обучение заменяется обученной
суррогатной моделью.

\begin{figure}[h]
    \centering
    \includegraphics[width=0.7\textwidth]{main_frame}
    \caption{Схема метода: суррогатная модель $u$ используется для подбора
    наиболее подходящих архитектур экспертов для каждого кластера данных.}
    \label{fig:workflow}
\end{figure}

Пусть обучающие данные разбиты на $M$ кластеров, и для
$\boldsymbol{b} \in \{0, 1\}^M$ пусть
$\mathcal{D}_{\boldsymbol{b}} = \bigcup_{m: b_m = 1} \mathcal{C}_m$ ---
подмножество кластеров, выбранное вектором $\boldsymbol{b}$. Суррогатная
функция $u(\boldsymbol{\alpha}, \boldsymbol{b})$ предсказывает среднюю
по объектам кросс-энтропийную ошибку (отрицательное среднее
лог-правдоподобие верного класса) архитектуры $\boldsymbol{\alpha}$,
обученной на $\mathcal{D}_{\boldsymbol{b}}$, без проведения фактического
обучения; это та же величина, что входит в
объектив~\eqref{eq:moe_objective} через
$e^{-u(\boldsymbol{\alpha}_k, \mathbf{r}_k)}$.

Суррогат обучается на датасете пар «архитектура--подмножество»,
собранных реальным обучением,
$\mathcal{D}_{\text{surr}} = \{(\boldsymbol{\alpha}^{(i)}, \boldsymbol{b}^{(i)}, y^{(i)})\}_{i=1}^{N_s}$,
где $y^{(i)}$ --- измеренная валидационная кросс-энтропия. Модель
$u_{\boldsymbol{\theta}}$ параметризуется графовой нейронной сетью (GNN),
кодирующей граф архитектуры, конкатенированной с эмбеддингом кластерной
маски $\boldsymbol{b}$; минимизируется среднеквадратичная ошибка
\begin{equation}
\label{eq:surr_training}
\min_{\boldsymbol{\theta}} \;
\frac{1}{|\mathcal{D}_{\text{surr}}|}
\sum_{(\boldsymbol{\alpha}, \boldsymbol{b}, y) \in \mathcal{D}_{\text{surr}}}
\bigl(u_{\boldsymbol{\theta}}(\boldsymbol{\alpha}, \boldsymbol{b}) - y\bigr)^2.
\end{equation}

Суррогат уточняется в цикле активного обучения. На каждой итерации с
помощью Monte-Carlo dropout получают предсказательное среднее
$\mu(\boldsymbol{\alpha}, \boldsymbol{b})$ и стандартное отклонение
$\sigma(\boldsymbol{\alpha}, \boldsymbol{b})$ по $T$ стохастическим
прогонам, после чего кандидаты оцениваются по нижней доверительной
границе $\mathrm{LCB} = \mu - \sigma$ (поскольку $u$ --- ошибка,
перспективны кандидаты с наименьшим предсказанным значением и высокой
эпистемической неопределённостью). Наиболее перспективный кандидат
обучается с нуля, полученная тройка добавляется в
$\mathcal{D}_{\text{surr}}$, и суррогат переобучается. Существенно, что
$u_{\boldsymbol{\theta}}$ обучается \emph{только на дискретных} масках
$\boldsymbol{b}$, поэтому дифференцируемая оптимизация должна сочетаться
с релаксацией дискретных входов.

\section{Суррогатно-ассистированный обобщённый EM (SGEM)}
\label{sec:sgem}

Рассмотрение маршрутизации «кластер--эксперт» как латентной переменной
превращает объектив~\eqref{eq:moe_objective} в модель смеси, к которой
напрямую применим EM-алгоритм, причём суррогат $u^{(t)}$ заменяет
труднодоступное истинное качество эксперта.

\textbf{Латентная форма целевой функции.} Для каждого объекта $n$ введём
унитарный (one-hot) латентный вектор $\mathbf{z}_n \in \{0,1\}^K$, где
$z_{nk} = 1$ тогда и только тогда, когда объект $n$ порождён экспертом
$k$, и соберём их в матрицу $\mathbf{Z} = \|z_{nk}\|$. Порождающая
история стандартна для смеси: эксперт объекта $n$ выбирается из
маршрутизации его кластера, $p(z_{nk} = 1) = r_{m(n), k}$, а метка
порождается с пообъектным правдоподобием эксперта
$e^{-u(\boldsymbol{\alpha}_k, \mathbf{r}_k)}$. Тогда правдоподобие
полных данных есть
\begin{equation}
\label{eq:complete_data}
p\bigl(\mathbf{y}, \mathbf{Z} \mid \boldsymbol{\alpha}_{1:K}, \mathbf{R}\bigr)
\;=\; \prod_{n=1}^{N} \prod_{k=1}^{K}
\Bigl( r_{m(n), k}\, e^{-u(\boldsymbol{\alpha}_k, \mathbf{r}_k)} \Bigr)^{z_{nk}},
\end{equation}
где роль весов смеси играет маршрутизация $r_{mk}$, а роль пообъектного
правдоподобия компоненты --- $e^{-u(\boldsymbol{\alpha}_k, \mathbf{r}_k)}$.

\emph{Эквивалентность с~\eqref{eq:moe_objective}.} Каждая строка
$\mathbf{z}_n$ унитарна, поэтому маргинализация сворачивает внутреннее
произведение в сумму; так как маршрутизация и суррогат постоянны внутри
кластера, группировка $|\mathcal{C}_m|$ одинаковых множителей даёт
$p(\mathbf{y} \mid \boldsymbol{\alpha}_{1:K}, \mathbf{R})
= \prod_{m} \bigl(\sum_{k} r_{mk}\, e^{-u(\boldsymbol{\alpha}_k, \mathbf{r}_k)}\bigr)^{|\mathcal{C}_m|}$,
логарифм чего в точности равен
$\mathcal{J}(\boldsymbol{\alpha}_{1:K}, \mathbf{R})$
из~\eqref{eq:moe_objective}. Таким образом, максимизация
объектива~\eqref{eq:moe_objective} есть оценка максимального
правдоподобия для латентной модели~\eqref{eq:complete_data}, которую мы
решаем EM-алгоритмом.

\textbf{E-шаг.} При фиксированных $(\boldsymbol{\alpha}_{1:K}, \mathbf{R})$
апостериорное распределение факторизуется по объектам, а ответственность
эксперта $k$ за кластер $m$ есть нормированный вес смеси
\begin{equation}
\label{eq:estep}
q_{mk} \;=\; \frac{r_{mk}\, e^{-u(\boldsymbol{\alpha}_k, \mathbf{r}_k)}}
{\sum_{j=1}^{K} r_{mj}\, e^{-u(\boldsymbol{\alpha}_j, \mathbf{r}_j)}},
\qquad \sum_{k=1}^{K} q_{mk} = 1.
\end{equation}
Интуитивно, кластер мягко приписывается к экспертам, которые и
претендуют на него ($r_{mk}$ велико), и предсказанно хорошо его
моделируют ($u(\boldsymbol{\alpha}_k, \mathbf{r}_k)$ мало).

\textbf{M-шаг.} Подстановка апостериора в математическое ожидание
лог-правдоподобия полных данных и группировка по кластерам дают
$Q$-функцию
\begin{equation}
\label{eq:qfun}
Q(\boldsymbol{\alpha}_{1:K}, \mathbf{R})
\;=\; \sum_{m=1}^{M} |\mathcal{C}_m| \sum_{k=1}^{K} q_{mk}
\bigl[\log r_{mk} \,-\, u(\boldsymbol{\alpha}_k, \mathbf{r}_k)\bigr],
\end{equation}
которая максимизируется по маршрутизации $\mathbf{R}$ и архитектурам
$\boldsymbol{\alpha}_{1:K}$. Два блока обновляются поочерёдно.

\emph{M-шаг по маршрутизации.} При фиксированных архитектурах $Q$
максимизируется по $\mathbf{R}$ (при ограничении
$\sum_k r_{mk} = 1$) градиентным подъёмом по логитам маршрутизации с
использованием straight-through Gumbel-Softmax для проброса градиента
через дискретные маски $\mathbf{r}_k$, которые потребляет суррогат.
Слагаемое $\sum_{m,k} |\mathcal{C}_m|\, q_{mk}\log r_{mk}$ есть
кросс-энтропия, чей условный оптимум $r_{mk} = q_{mk}$ подтягивает
маршрутизацию к текущим ответственностям, а слагаемое
$-q_{mk}\,u(\boldsymbol{\alpha}_k, \mathbf{r}_k)$ дополнительно смещает
столбец $\mathbf{r}_k$ к маскам, на которых эксперт предсказанно хорош.

\emph{M-шаг по архитектурам.} При фиксированной маршрутизации
зависящая от архитектуры часть~\eqref{eq:qfun} есть
$-\sum_{m} |\mathcal{C}_m|\, q_{mk}\, u(\boldsymbol{\alpha}_k, \mathbf{r}_k)$,
поэтому каждый эксперт оптимизируется независимо \emph{минимизацией}
своей взвешенной ответственностями суррогатной ошибки. Так как $u$
недифференцируема по $\boldsymbol{\alpha}$, кандидаты-архитектуры
сэмплируются из пространства поиска, и для каждого эксперта $k$
оставляется кандидат с наименьшим $u(\boldsymbol{\alpha}_k, \mathbf{r}_k)$.

Алгоритм SGEM суммируется следующей процедурой. Вход: суррогат
$u_{\boldsymbol{\theta}}$, пространство поиска $\mathcal{A}$, число
экспертов $K$, число кластеров $M$, число итераций $T$.
\begin{enumerate}
    \item Инициализировать архитектуры
    $\boldsymbol{\alpha}_1, \dots, \boldsymbol{\alpha}_K$ случайно и
    логиты маршрутизации
    $\boldsymbol{\ell} \in \mathbb{R}^{M\times K} \sim \mathcal{N}(\mathbf{0}, I)$.
    \item Для $t = 1, \ldots, T$ повторять:
    \begin{enumerate}
        \item \emph{E-шаг}: вычислить ответственности
        $q_{mk}$ по~\eqref{eq:estep};
        \item \emph{M-шаг по $\mathbf{R}$}: градиентный подъём
        по~\eqref{eq:qfun} с Gumbel-Softmax;
        \item \emph{M-шаг по $\boldsymbol{\alpha}$}: для каждого эксперта
        сэмплировать кандидатов и оставить с наименьшим
        $u(\cdot, \mathbf{r}_k)$;
        \item \emph{Шаг активного обучения}: обучить выбранных кандидатов,
        пополнить $\mathcal{D}_{\text{surr}}$, переобучить
        $u_{\boldsymbol{\theta}}$.
    \end{enumerate}
    \item Вернуть $\boldsymbol{\alpha}_{1:K}$ и жёсткое назначение
    $\hat{r}_{mk} = [\,k = \arg\max_{k'} r_{mk'}\,]$.
\end{enumerate}

\section{Анализ сходимости}
\label{ch:theory}

Особенность Алгоритма SGEM в том, что суррогат $u^{(t)}$, используемый на
E- и M-шагах, сам обновляется между итерациями, поэтому классический
аргумент монотонности EM~\cite{dempster1977em, wu1983convergence}
напрямую неприменим. Покажем, что при достаточно быстром убывании ошибки
суррогата SGEM сходится в том же смысле, что и обобщённый EM.

Пусть $u_{\mathrm{true}}$ --- неизвестная истинная ошибка эксперта
(пообъектная кросс-энтропия архитектуры $\boldsymbol{\alpha}$, обученной
на подмножестве $\mathbf{r}$). Определим пошаговую ошибку суррогата в
шкале потерь
\begin{equation}
\label{eq:eps_t}
\varepsilon_t \;=\; \sup_{\boldsymbol{\alpha} \in \mathcal{A}, \;
\mathbf{r} \in [0, 1]^M}
\bigl|u^{(t)}(\boldsymbol{\alpha}, \mathbf{r})
- u_{\mathrm{true}}(\boldsymbol{\alpha}, \mathbf{r})\bigr|.
\end{equation}
Параметр $\theta = (\boldsymbol{\alpha}_{1:K}, \mathbf{R})$ лежит в
компактном множестве $\Theta = \mathcal{A}^K \times (\Delta_K)^M$;
обозначим $N = \sum_m |\mathcal{C}_m|$. Истинное (взвешенное по размерам
кластеров) лог-правдоподобие есть
\begin{equation}
\label{eq:Ltrue}
L_{\mathrm{true}}(\theta) \;=\; \sum_{m=1}^M |\mathcal{C}_m|\,\log
\sum_{k=1}^K r_{mk}\, e^{-u_{\mathrm{true}}(\boldsymbol{\alpha}_k, \mathbf{r}_k)},
\end{equation}
а истинная $Q$-функция ---
\begin{equation}
\label{eq:Qtrue}
Q_{\mathrm{true}}(\theta; \theta') \;=\; \sum_{m=1}^M |\mathcal{C}_m|
\sum_{k=1}^K w_{mk}(\theta')
\bigl[\log r_{mk} - u_{\mathrm{true}}(\boldsymbol{\alpha}_k, \mathbf{r}_k)\bigr],
\end{equation}
где $w_{mk}(\theta')$ --- ответственности, вычисленные при
$u_{\mathrm{true}}$. Введём вариационную свободную энергию,
параметризованную функцией потерь $u$:
\[
\mathcal{F}_u(q, \theta) \;=\; \sum_{m=1}^{M} |\mathcal{C}_m|
\sum_{k=1}^{K} q_{mk}\, \log \frac{r_{mk}\, e^{-u(\boldsymbol{\alpha}_k, \mathbf{r}_k)}}{q_{mk}},
\qquad L_u(\theta) = \max_q \mathcal{F}_u(q, \theta).
\]

\begin{assumption}
\label{ass:bounded}
Существует $U_{\max} < \infty$ такое, что
$0 \leq u_{\mathrm{true}}, u^{(t)} \leq U_{\max}$ для всех
$\boldsymbol{\alpha}, \mathbf{r}, t$.
\end{assumption}

\begin{assumption}
\label{ass:summable}
Ошибки суррогата суммируемы:
$\sum_{t=0}^{\infty} \varepsilon_t < \infty$.
\end{assumption}

\begin{assumption}
\label{ass:gem}
M-шаг обобщён, т.е. для каждого $t$ возвращается $\theta^{(t+1)}$ с
$\mathcal{F}_{u^{(t)}}(q^{(t)}, \theta^{(t+1)})
\geq \mathcal{F}_{u^{(t)}}(q^{(t)}, \theta^{(t)})$.
\end{assumption}

\begin{assumption}
\label{ass:cont}
Для каждого $\boldsymbol{\alpha}$ отображение
$\mathbf{r}\mapsto u_{\mathrm{true}}(\boldsymbol{\alpha}, \mathbf{r})$
непрерывно на $[0,1]^M$.
\end{assumption}

\begin{theorem}[Сходимость SGEM с адаптивным суррогатом]
\label{thm:sgem}
При выполнении Предположений~\ref{ass:bounded}--\ref{ass:cont}
последовательность $\{\theta^{(t)}\}$, порождаемая Алгоритмом SGEM,
удовлетворяет:
\begin{enumerate}
\item[\textup{(a)}] последовательность
$\{L_{\mathrm{true}}(\theta^{(t)})\}$ сходится;
\item[\textup{(b)}] всякая предельная точка $\theta^{\star}$ является
стационарной точкой $L_{\mathrm{true}}$ на $\Theta$ в смысле
$Q_{\mathrm{true}}(\theta;\theta^{\star}) \leq
Q_{\mathrm{true}}(\theta^{\star};\theta^{\star})$ для всех
$\theta\in\Theta$.
\end{enumerate}
\end{theorem}

\begin{proof}[Схема доказательства]
Ключевые шаги следующие. Из вариационного представления
$L_u(\theta) = \max_q \mathcal{F}_u(q, \theta)$ с максимумом в точке
$q^{\star}_{mk} \propto r_{mk} e^{-u(\boldsymbol{\alpha}_k, \mathbf{r}_k)}$
следует равномерная оценка отклонения:
$|\mathcal{F}_{u_{\mathrm{true}}}(q,\theta) - \mathcal{F}_{u^{(t)}}(q,\theta)|
\leq N\varepsilon_t$ и, как следствие,
$|L_{\mathrm{true}}(\theta) - L^{(t)}(\theta)| \leq N\varepsilon_t$.
Объединяя это с Предположением~\ref{ass:gem} (подъём свободной энергии на
M-шаге) и равенством на E-шаге, получаем
\emph{квазимонотонность}
\begin{equation}
\label{eq:quasi-mono}
L_{\mathrm{true}}(\theta^{(t+1)}) \;\geq\;
L_{\mathrm{true}}(\theta^{(t)}) \;-\; 2N\,\varepsilon_t.
\end{equation}
Положив $V_t = L_{\max} - L_{\mathrm{true}}(\theta^{(t)}) \geq 0$
(величина $L_{\max} = \sup_\theta L_{\mathrm{true}}$ конечна ввиду
компактности $\Theta$), из~\eqref{eq:quasi-mono} имеем
$V_{t+1} \leq V_t + 2N\varepsilon_t$ при
$\sum_t 2N\varepsilon_t < \infty$. По теореме Роббинса--Зигмунда
\cite{robbins1971convergence} последовательность $\{V_t\}$ сходится, что
доказывает~(a). Для~(b) используется устойчивость softmax-ответственностей
к возмущению суррогата: при $\varepsilon_t \to 0$ ответственности и
$Q$-функция, вычисленные при $u^{(t)}$, сходятся к истинным; если бы
предельная точка не была стационарной, нашёлся бы $\theta'$ с
$Q_{\mathrm{true}}(\theta';\theta^{\star}) -
Q_{\mathrm{true}}(\theta^{\star};\theta^{\star}) = \delta > 0$, что по
переходу через $\tilde Q^{(t)}$ и подъёму на M-шаге дало бы прирост
$\Delta_t \geq \delta/8$ для бесконечного числа $t$, противоречащий
суммируемости из~(a). Следовательно, $\theta^{\star}$ стационарна.
\end{proof}

Побочным следствием доказательства пункта~(a) является явная оценка
скорости
\[
L_{\mathrm{true}}(\theta^{\star}) - L_{\mathrm{true}}(\theta^{(T)})
\;\leq\; 2N \sum_{t = T}^{\infty} \varepsilon_t.
\]
Если $\varepsilon_t = O(t^{-p})$ при $p > 1$, то
$L_{\mathrm{true}}(\theta^{(T)}) \to L_{\mathrm{true}}(\theta^{\star})$
со скоростью $O(T^{1 - p})$, что согласуется со скоростями,
наблюдаемыми в анализе стохастического EM~\cite{cappe2009online}.

\begin{remark}
Предположение суммируемости~\ref{ass:summable} --- суррогатный аналог
условия убывающего шага в стохастическом EM; ослабление его до
$\varepsilon_t \to 0$ сохраняет пункт~(b), но не пункт~(a). На практике
суммируемость достигается, когда шаг активного обучения обеспечивает
достаточно полное покрытие пространства
$\mathcal{A}\times[0,1]^M$.
\end{remark}

\textbf{Новизна и связь с классическими результатами о сходимости.}
Классическая теория сходимости EM и его обобщённого варианта
\cite{dempster1977em, wu1983convergence} опирается на монотонный подъём
единого \emph{фиксированного} правдоподобия: каждый M-шаг не уменьшает
одну и ту же целевую функцию, и сходимость следует из монотонной
ограниченной последовательности. В нашей постановке эта предпосылка
нарушена. Поскольку качество эксперта берётся с адаптивно
переобучаемого суррогата $u^{(t)}$, а не из точного обучения,
оптимизируемая на итерации $t$ цель сама является «движущейся мишенью»,
а истинное правдоподобие $L_{\mathrm{true}}$, которое нас в итоге
интересует, алгоритмом \emph{никогда} не вычисляется. Поэтому точная
монотонность недоступна, и стандартный аргумент неприменим.

Теорема~\ref{thm:sgem} показывает, что сходимость тем не менее
сохраняется, --- в этом состоит основной теоретический вклад. Механизм
--- \emph{квази}-монотонность
$L_{\mathrm{true}}(\theta^{(t+1)}) \geq L_{\mathrm{true}}(\theta^{(t)}) - 2N\varepsilon_t$:
истинное правдоподобие может проседать, но лишь на величину,
пропорциональную пошаговой ошибке суррогата $\varepsilon_t$.
Суммируемость этих ошибок превращает
$\{L_{\mathrm{true}}(\theta^{(t)})\}$ в почти-супермартингал, и аргумент
Роббинса--Зигмунда~\cite{robbins1971convergence} даёт и сходимость, и
стационарность всякой предельной точки \emph{истинной} цели --- хотя
каждый E- и M-шаг видит только суррогат.

Это отличается от анализа стохастического и онлайн-EM
\cite{cappe2009online}, где возмущение точного шага EM ---
\emph{нулевое в среднем стохастическое} (от Монте-Карло E-шага),
гасимое убывающим шагом. В нашем случае возмущение --- наоборот,
\emph{детерминированная ограниченная ошибка аппроксимации} обученного
суррогата, а роль убывающего шага играет условие
суммируемости~\ref{ass:summable}, малость которого обеспечивается не
усреднением шума, а покрытием пространства поиска при активном обучении.
Насколько нам известно, это первая гарантия сходимости для EM-подобной
процедуры архитектурного поиска, чья пошаговая цель задаётся
\emph{адаптивно обучаемой} суррогатной моделью.

